\section{Packing from a vertex cover}\label{sec:completeness}

\begin{proposition}\label{prop:completeness}
If the input graph has a vertex cover of size at most $k$, the constructed
items can be packed into $B$ bins.
\end{proposition}

\begin{proof}
Enlarge the cover to a set $C\subseteq V$ of exactly $k$ vertices.
Call vertices of $C$ \emph{selected} and the others \emph{unselected}.
We first construct the $\mathrm M$ bins separately at each vertex.
The local construction will also be useful when no common cover is
specified.

\paragraph{Rows and auxiliary slots.}
At label $v$, put every $\Tp$ and $\Tm$ row into $\mathrm M$.
At each position in $\mathcal R_v$, put the $d$ long $\Srow$ rows
into $\mathrm M$ if $v$ is selected, and the $d$ short rows otherwise.
There are $t_v+J_v$ such rows, hence this many slots for each of
$\Zanc$ and $\DM$.
The roots and padding give the useful bound
\begin{equation}\label{eq:local-slot-bound}
 t_v=P+d\bigl(|V(T^+)|+|V(T^-)|\bigr)\ge P+2d.
\end{equation}
In particular the number of zero-length $\DM$ copies is
$t_v+J_v-P\ge2d+J_v$.

At a selected vertex, choose $d$ distinct plus-root rows and give
them length-one $\Up$ items; at an unselected vertex do the same
with minus-root rows and length-one $\Um$ items.  Reserve $d$
zero-length $\DM$ copies for these funded rows.
Every other $\mathrm M$ row receives a length-zero global auxiliary:
$\Tp$ and $\Srow$ require $\Up$, and $\Tm$ requires $\Um$.
Thus label $v$ has $t_{v,+}+J_v$ plus slots and $t_{v,-}$ minus
slots.  Globally the requested length-one counts are exactly $dk$
and $d(n-k)$; the total slots are precisely the respective global
inventories.  Consequently all global auxiliaries of these two
subclasses can be assigned as distinct physical copies, exhausting
their length-one and length-zero stocks.

In the tree on the unfunded side, choose the separate $d$-fold
marking from Lemma~\ref{lem:competing-trees}.  For every marked
depth-$\ell$ vertex, choose one of its $d$ endpoint rows and attach
a $\DM$ item of length $\lambda_\ell$.  The depth quota supplies
enough such items.  These marked rows are distinct and have positive
depth, so none is a funded-root row.  If $s_v$ nodes are marked, then
$s_v\le P$.  After these assignments and the reserved root zeros,
there are $t_v+J_v-s_v-d$ unassigned local slots.  By
Equation~\eqref{eq:local-slot-bound}, this is at least $P-s_v$, the
number of unused positive local items.  Put those items in distinct
remaining slots, then fill all other slots with the remaining zeros.
This uses every $\DM$ copy.  All additional positive lengths move
their rows earlier and cannot remove any coverage already obtained.
The same argument applies when $v$ is isolated and $J_v=0$.

Attach a $\FT$ flag to every tree row and a $\FM$ flag to every
$\Srow$ row used here.  These requests exhaust the $\FT$ and $\FM$ inventories
after summing over labels.  From now on keep each row, its two
auxiliaries, and its flag together.

\paragraph{Coverage and deadlines.}
For each such group, let $b$ be its row baseline and let
\[
 h=b-\Delta\ind\{\text{row is short}\}
      -\len(U)-\len(D)
\]
be its completion coordinate, as in
Lemma~\ref{lem:good-feasibility}.  Always $h\le b$.
Each funded-root row covers $[0,1)$ or $[4,5)$ through its interval
$[h,b)$.  Hence the $d$ funded roots cover every test interval on
that side $d$ times, including all job intervals when plus is funded.
At a marked depth-$\ell$ node on the other side, the assigned local
length gives precisely its cell with the right endpoint removed.
By nesting in Lemma~\ref{lem:geometry}, this interval contains the
test interval of every descendant leaf.  The $d$-fold marking covers every leaf test on that
side at least $d$ times.  Finally, at an unselected vertex each job
has $d$ short rows in $\mathrm M$; even before any auxiliary
shortening they cover its interval $[b_r-\Delta,b_r)$.
Thus every test interval receives at least $d$ covering intervals
in either state.

The baseline multiset of these $\mathrm M$ rows is exactly the one
defining $B_v(a)$: all tree baselines and $d$ copies of $b_r$ at
each job, irrespective of the choice of short or long rows.
Recall that $\mathcal I_v$ is the union of the test intervals at $v$.
For every real $a$,
\begin{equation}\label{eq:complete-prefix}
 \#\{h\le a\}
 =B_v(a)+\#\{h\le a<b\}
 \ge B_v(a)+d\ind\{a\in\mathcal I_v\}.
\end{equation}
The last expression is the exact number of deadlines at most $a$.
Here the left-closed, right-open convention includes a shifted
deadline at the left endpoint and stops counting the extra copy
when its original baseline enters at the right endpoint.

Sort the completions as $h_1\le\cdots\le h_{t_v+J_v}$ and the
deadlines as $z_1\le\cdots\le z_{t_v+J_v}$, retaining the identity
of each copy.  Equation~\eqref{eq:complete-prefix} at $a=z_j$
shows $h_j\le z_j$.  Attach the $j$th deadline anchor to the group
with completion $h_j$.
Only the $\Zanc$ anchors are permuted; the row, global auxiliary,
local auxiliary, and flag stay together.  All the anchors have the
same subclass and label $v$, so every resulting tuple keeps its
required pattern and equal labels.  Its deadline inequality makes
it feasible by Lemma~\ref{lem:good-feasibility}.

\paragraph{The remaining rows and edge resources.}
At each incidence position, pair its remaining $d$ rows with its
$d$ keys, using the exact position and label.  Add distinct $\DG$
copies of that label and $\FG$ flags.  The remaining rows are short
at a selected vertex and long at an unselected vertex.
For a $\Waux$ item of the exact edge, the coordinate condition in
Lemma~\ref{lem:good-feasibility} reduces to
\begin{equation}\label{eq:exact-edge-completeness}
 \Delta\bigl(\ind\{\text{nonpermit}\}
                 -\ind\{\text{row is short}\}\bigr)\le0.
\end{equation}
Thus a long row requires a permit, whereas a short row accepts
either species.

Each endpoint of an edge has $dR$ such slots.  If exactly one
endpoint is unselected, assign all $dR$ permits to its long rows
and all $dR$ nonpermits to the selected endpoint's short rows.
If both endpoints are selected, all $2dR$ rows are short; assign
the $dR$ permits and the $dR$ nonpermits to these slots in any order.
Since $C$ is a cover, these are the only cases.  In both cases
every $\Waux$ copy is used, and all $\mathrm G$ tuples are feasible.

At label $v$ we have made $t_v+J_v$ bins of type $\mathrm M$ and
$J_v$ of type $\mathrm G$.  They use every labeled row, anchor,
and local auxiliary once.  Their flag counts are $t_v$ of $\FT$
and $J_v$ each of $\FM$ and $\FG$, and all global inventories have
been exhausted.  Every bin follows Table~\ref{tab:patterns} and
has equal labels, so Lemmas~\ref{lem:encoding} and
\ref{lem:good-feasibility} apply to the actual encoded sizes.
The construction is a partition of all items into
$\sum_v(t_v+2J_v)=B$ feasible bins.
\end{proof}

\section{The fractional packing}\label{sec:fractional}

We construct a fractional packing of value $B$ for the four-vertex
instance. Both configuration models are those of
Definition~\ref{def:configuration-lp}; all feasible configurations are
allowed, with no restriction to the table patterns used in our witness.
The matching lower bound will follow from the five-item capacity limit.

\begin{proposition}\label{prop:fractional}
For the instance constructed from the complete graph on four
vertices with $k=2$, using the parameters in
Equation~\eqref{eq:parameters}, both configuration relaxations have
value $B$:
\[
 \OPT_{\mathrm{LP,ind}}=\LP=B.
\]
\end{proposition}

\begin{proof}
For each label $v$, construct the local groups from the proof of
Proposition~\ref{prop:completeness} twice, once selected and once
unselected.  In the selected state, put short rows in $\mathrm G$
and require only nonpermits of the exact edge.  In the unselected
state, put long rows there and require only permits of the exact
edge.  Equation~\eqref{eq:exact-edge-completeness} holds with
equality in both states.
The local $\mathrm M$ construction and deadline matching did not
use neighboring states.  Each state therefore gives a collection
of $N_v=t_v+2J_v$ feasible bin templates, using every physical
copy in the labeled roles $X,A,D$ exactly once.
Fix these labeled assignments in both states, including the local
auxiliaries and the independently sorted anchors.  The two remaining
slots in a template specify a global-auxiliary species and a flag
species.

\paragraph{Exact resource demands.}
Table~\ref{tab:local-state-demands} lists the required counts.
The two $\Waux$ rows refer to each edge $e$ incident to $v$.
\begin{table}[htbp]
\centering
\caption{Pooled-item demands at one vertex in the two local states.}
\label{tab:local-state-demands}
\vspace{6pt}
\begin{tabular}{@{}lccc@{}}
\toprule
Species & Selected & Unselected & Half of each\\
\midrule
$\Up$, length $1$ & $d$ & $0$ & $d/2$\\
$\Up$, length $0$ & $t_{v,+}+J_v-d$ & $t_{v,+}+J_v$
 & $t_{v,+}+J_v-d/2$\\
$\Um$, length $1$ & $0$ & $d$ & $d/2$\\
$\Um$, length $0$ & $t_{v,-}$ & $t_{v,-}-d$ & $t_{v,-}-d/2$\\
$\Waux$, permit of $e$ & $0$ & $dR$ & $dR/2$\\
$\Waux$, nonpermit of $e$ & $dR$ & $0$ & $dR/2$\\
$\FT$ & $t_v$ & $t_v$ & $t_v$\\
$\FM$ & $J_v$ & $J_v$ & $J_v$\\
$\FG$ & $J_v$ & $J_v$ & $J_v$\\
\bottomrule
\end{tabular}
\end{table}

Each entire local state can in fact choose distinct pooled copies
from the actual inventories.  Both length-one pools have $2d$
copies, while a state needs at most $d$ of either.
Every other vertex contributes at least $d$ root slots on each
side, so the zero-length stocks satisfy
\begin{align*}
 \sum_w(t_{w,+}+J_w)-2d&\ge t_{v,+}+J_v+d,\\
 \sum_w t_{w,-}-2d&\ge t_{v,-}+d.
\end{align*}
These bounds exceed the respective largest local demands.
Each permit or nonpermit pool of an incident edge has exactly
$dR$ copies, enough for its local demand, and each flag pool
contains at least the number required at one vertex.  All pools
requested by any template are therefore nonempty.  Different
local states may reuse copies; we will now allocate their
fractional uses explicitly.

Taking half of both states at every vertex gives total demand equal
to every pooled stock.
Indeed, summing over the four vertices gives $2d$ length-one
items on each side and exactly the prescribed zero-length counts.
Each edge has two endpoints, which together request $dR$ permits
and $dR$ nonpermits.  Each flag has the same demand in both states,
so its weighted total is also its inventory.

\paragraph{Columns using physical copies.}
For a template $b$, let $i_b^X,i_b^A,i_b^D$ be its fixed labeled
copies.  Let $\mathcal U_b$ be the set of physical copies in its
required global species, distinguishing side and length for
$\Up,\Um$, and edge and permit status for $\Waux$.
Let $\mathcal F_b$ be the set of copies of its required flag.
For every pair $(u,f)\in\mathcal U_b\times\mathcal F_b$, assign
the individual-copy configuration
\begin{equation}\label{eq:physical-fractional-columns}
 H_{b,u,f}=\{i_b^X,i_b^A,u,i_b^D,f\}
 \quad\text{weight}\quad
 \frac{1}{2|\mathcal U_b|\,|\mathcal F_b|}.
\end{equation}
Do this for every template in both states of every vertex, adding
weights whenever the same configuration is generated more than once.
All five physical copies in a column are distinct because their
roles are disjoint.  Copies within either requested pool have
identical encoded sizes; hence every generated column is feasible.
Coincidences of numerical sizes between different declared species
do not affect their disjoint physical indices or this argument.

The columns generated by one template have total weight $1/2$.
Every labeled item occurs in one template in each of its two local
states, giving total coverage $1$.
For a global pool $\mathcal U$, each requesting template contributes
$1/(2|\mathcal U|)$ to each copy's coverage after summing over flag
choices.  Its total half-weight demand equals $|\mathcal U|$, as
just proved, so each global copy also has coverage $1$.
The same calculation, summing over global choices, gives coverage
$1$ for each flag copy.  Equation~\eqref{eq:physical-fractional-columns}
therefore defines a feasible solution of
Equation~\eqref{eq:individual-configuration-lp}, with objective
\[
 \sum_{v=1}^4\left(\frac{N_v}{2}+\frac{N_v}{2}\right)=B.
\]

Project each individual configuration to its vector of multiplicities
of equal numerical sizes, and sum the weights of configurations
with the same vector.  Capacity is unchanged, and summing the
individual coverage equalities over copies of a given size gives
its required multiplicity.  This proves $\LP\le B$ for the ordinary
type LP, including when different species have equal sizes.

Finally, Lemma~\ref{lem:encoding} gives $a_i>1/6$ for every item.
Every feasible configuration consequently has cardinality at most
five, counting multiplicity in the type formulation.  Giving dual
weight $1/5$ to every physical item, or to every numerical size
type, is therefore feasible for the respective covering LP.
Both dual objectives equal $|\mathcal I|/5=5B/5=B$.
Thus both relaxations have value exactly $B$.

This construction uses feasible local bin configurations.  It does
not require a vertex cover of size two in the four-vertex clique,
nor an integral packing obtained by making those local choices
simultaneously.  The integral lower bound will follow separately
from Proposition~\ref{prop:soundness}.
\end{proof}
